\documentclass[11pt,a4paper]{article}

\usepackage[utf8]{inputenc}
\usepackage[T1]{fontenc}
\usepackage[french]{babel}
\usepackage{amsmath,amssymb,amsthm}
\usepackage{geometry}
\usepackage{hyperref}
\geometry{margin=2.5cm}

\newcommand{\dd}{\mathrm{d}}
\newcommand{\R}{\mathbb{R}}
\newcommand{\indic}[1]{\mathrm{1}_{#1}}

\title{Formalisme jump-diffusion pour la contraction musculaire \\
\large Équation maîtresse et fermeture à force imposée}
\author{}
\date{}

\begin{document}

\maketitle

\section{Contexte}

On considère le modèle jump-diffusion de l'interaction actine--myosine
proposé par Chaintron, Kimmig, Caruel et Moireau
(\emph{J. Appl. Phys.} \textbf{134}, 194901 (2023),
\href{https://doi.org/10.1063/5.0158191}{doi:10.1063/5.0158191}).
L'état d'une tête de myosine représentative est décrit par le processus
de Markov
\[
Z_t = (\alpha_t, X_t, Y_t) \in \{0,1\} \times \left(-\tfrac{d}{2}, +\tfrac{d}{2}\right) \times \R,
\]
où $\alpha_t$ est l'état d'attachement (0 : détaché, 1 : attaché),
$X_t$ la position de la tête par rapport à son point d'ancrage, et
$Y_t$ la variable de conformation associée au \emph{power-stroke}.

\section{Équation maîtresse du modèle (éq. 17 du papier)}

Le triplet $(\alpha_t, X_t, Y_t)$ est régi par le système d'équations
différentielles stochastiques suivant.

\subsection*{Dynamique de $\alpha_t$}

\begin{align*}
\dd\alpha_t
&= \indic{\alpha_t=0}\left[
\int_{\R_+} \indic{z \le K_{0\to1}(X_t,Y_t,s(t))}\, N_{0\to1}(\dd t,\dd z)
+ \int_{\R_+} \indic{z \le K^{\mathrm{rev}}_{1\to0}(X_t,Y_t,s(t))}\, N^{\mathrm{rev}}_{1\to0}(\dd t,\dd z)
\right] \\[1mm]
&\quad - \indic{\alpha_t=1}\left[
\int_{-d/2}^{d/2}\!\!\int_{\R_+} \indic{z \le K_{1\to0}(x,Y_t,s(t))}\, N_{1\to0}(\dd t,\dd x,\dd z) \right.\\
&\qquad\qquad\quad \left. + \int_{-d/2}^{d/2}\!\!\int_{\R_+} \indic{z \le K^{\mathrm{rev}}_{0\to1}(x,Y_t,s(t))}\, N^{\mathrm{rev}}_{0\to1}(\dd t,\dd x,\dd z)
\right]
\end{align*}

\subsection*{Dynamique de $X_t$}

\begin{align*}
\dd X_t
&= \indic{\alpha_t=0}\left[-\eta_x^{-1}\partial_x w_0(X_t,Y_t)\,\dd t
+ \sqrt{2\eta_x^{-1}k_BT}\; \dd B_t^x\right]
+ \indic{\alpha_t=1}\, \dot x_c(t)\,\dd t \\[1mm]
&\quad + \indic{\alpha_t=0}\left[
\int_{\R_+} (s(t)-X_t)\,\indic{z \le K_{0\to1}(X_t,Y_t,s(t))}\, N_{0\to1}(\dd t,\dd z) \right.\\
&\qquad\qquad\quad \left. + \int_{\R_+} (s(t)-X_t)\,\indic{z \le K^{\mathrm{rev}}_{1\to0}(X_t,Y_t,s(t))}\, N^{\mathrm{rev}}_{1\to0}(\dd t,\dd z)
\right] \\[1mm]
&\quad + \indic{\alpha_t=1}\left[
\int_{-d/2}^{d/2}\!\!\int_{\R_+} (x-s(t))\,\indic{z \le K_{1\to0}(x,Y_t,s(t))}\, N_{1\to0}(\dd t,\dd x,\dd z) \right.\\
&\qquad\qquad\quad \left. + \int_{-d/2}^{d/2}\!\!\int_{\R_+} (x-X_t)\,\indic{z \le K^{\mathrm{rev}}_{0\to1}(x,Y_t,s(t))}\, N^{\mathrm{rev}}_{0\to1}(\dd t,\dd x,\dd z)
\right]
\end{align*}

\subsection*{Dynamique de $Y_t$}

\[
\dd Y_t = -\eta_y^{-1}\partial_y w_{\alpha_t}(X_t,Y_t)\,\dd t
+ \sqrt{2\eta_y^{-1}k_BT}\; \dd B_t^y
\]

\bigskip
\noindent
Ici, $N_{0\to1}$, $N^{\mathrm{rev}}_{0\to1}$, $N_{1\to0}$, $N^{\mathrm{rev}}_{1\to0}$
sont quatre mesures aléatoires de Poisson indépendantes associées
respectivement à l'attachement direct, au détachement inverse,
au détachement direct et à l'attachement inverse ; $B_t^x, B_t^y$
sont deux mouvements browniens indépendants (bruit thermique) ;
$\eta_x,\eta_y$ sont les coefficients de friction ; $k_B$ la constante
de Boltzmann, $T$ la température, et $w_\alpha(x,y)$ le paysage
énergétique de la tête dans l'état $\alpha$.

\section{Quantité dérivée $Z_t$ (somme sur la population)}

On définit désormais, pour une population de $N$ têtes indexées par
$i$, la quantité agrégée
\[
Z_t := \sum_i \alpha_i^t\,\kappa\,(X_i^t+Y_i^t),
\]
où $\kappa$ est la raideur de la myosine. Cette quantité intervient
dans la relation de fermeture à force imposée $F$ :
\[
\frac{\dd s}{\dd t} = \frac{F}{v} - \frac{1}{v}\,Z_t,
\]
avec $v$ un paramètre de viscosité.

\subsection*{Remarque}

Dans le modèle du papier, $\kappa(s+y) = \partial_s w_1(s,y)$
correspond exactement à la force générée par une tête attachée
(cf.\ définition de $\tau_c(t)$, section III.C du papier). La quantité
$Z_t/N$ correspond donc à la force macroscopique moyenne
$\tau_c(t)$, ce qui assure la cohérence de la fermeture proposée
avec le cadre thermodynamique de l'article.

\section{Dérivation de la dynamique de $Z_t$}

\subsection{Dynamique d'une seule tête}

Soit $\xi_i^t := \alpha_i^t\,\kappa\,(X_i^t+Y_i^t) = \varphi(\alpha_i^t, X_i^t, Y_i^t)$
avec $\varphi(\alpha,x,y) = \alpha\,\kappa\,(x+y)$.

Comme $\varphi$ est \emph{affine} en $(x,y)$ à $\alpha$ fixé, les
termes du second ordre de la formule d'Itô (dérivées
$\partial^2_{xx}$, $\partial^2_{yy}$) s'annulent. Il reste :

\paragraph{Partie continue.} Tant que $\alpha_i^t=1$ (donc
$X_i^t=s(t)$) :
\[
\dd\xi_i^t\Big|_{\text{cont}}
= \indic{\alpha_i^t=1}\,\kappa\Big[\dot x_c(t)
- \eta_y^{-1}\partial_y w_1(X_i^t,Y_i^t)\Big]\dd t
+ \indic{\alpha_i^t=1}\,\kappa\sqrt{2\eta_y^{-1}k_BT}\;\dd B_i^{y,t}.
\]
Lorsque $\alpha_i^t=0$, on a $\xi_i^t \equiv 0$ : pas de contribution
continue.

\paragraph{Partie sauts.} À un attachement (direct ou inverse), la
position saute instantanément à $s(t)$, donc $\xi_i^t$ saute de $0$ à
$\kappa(s(t)+Y_i^t)$ ; à un détachement, l'inverse se produit :
\begin{align*}
\dd\xi_i^t\Big|_{\text{sauts}}
= \kappa\big(s(t)+Y_i^t\big)\Big[
&\indic{\alpha_i^t=0}\big(\dd N_{0\to1,i} + \dd N^{\mathrm{rev}}_{1\to0,i}\big) \\
- &\indic{\alpha_i^t=1}\big(\dd N_{1\to0,i} + \dd N^{\mathrm{rev}}_{0\to1,i}\big)
\Big],
\end{align*}
où $\dd N_{\cdot,i}$ désigne le comptage des sauts acceptés (intégrale
de $\indic{z\le K(\cdot)}$ contre la mesure de Poisson correspondante,
comme dans l'équation maîtresse pour $\dd\alpha_i^t$).

\subsection{Sommation sur la population}

En sommant sur toutes les têtes $i$ (les mesures de Poisson et les
browniens étant indépendants d'une tête à l'autre) :

\begin{align*}
\dd Z_t
&= \sum_{i:\,\alpha_i^t=1} \kappa\Big[\dot x_c(t)
- \eta_y^{-1}\partial_y w_1(X_i^t,Y_i^t)\Big]\dd t
+ \sum_{i:\,\alpha_i^t=1} \kappa\sqrt{2\eta_y^{-1}k_BT}\;\dd B_i^{y,t} \\[1mm]
&\quad + \sum_i \kappa\big(s(t)+Y_i^t\big)\Big[
\indic{\alpha_i^t=0}\big(\dd N_{0\to1,i} + \dd N^{\mathrm{rev}}_{1\to0,i}\big)
- \indic{\alpha_i^t=1}\big(\dd N_{1\to0,i} + \dd N^{\mathrm{rev}}_{0\to1,i}\big)
\Big].
\end{align*}

\subsection{Système fermé}

Avec la relation de fermeture
\[
\frac{\dd s}{\dd t} = \frac{F}{v} - \frac{Z_t}{v},
\]
le système devient auto-cohérent : $s(t)$ (et donc $\dot x_c(t) = \dot s(t)$,
paramètre exogène dans le modèle original à vitesse imposée) devient
une variable dynamique pilotée par $Z_t$, qui dépend en retour de
l'état de l'ensemble de la population de têtes attachées.

\bigskip
\noindent
\textbf{Remarque.} Cette relation de fermeture à force imposée $F$
n'est pas présente telle quelle dans le papier original ; elle en
constitue une extension naturelle, dans l'esprit de la perspective
évoquée par les auteurs en conclusion (simulation du modèle en
conditions de force imposée, avec prise en compte explicite de
l'élasticité du filament).

\end{document}